\chapter{A camada de agente}
\label{cap:agente}

\begin{objetivos}
Ao final deste capítulo, você será capaz de:
\begin{itemize}[leftmargin=*,itemsep=0.2em]
  \item descrever uma cadeia de pré-processamento que produza sinais comparáveis entre
        poços;
  \item calcular as treze métricas por sensor que traduzem um segmento anômalo em
        descrição quantitativa;
  \item descrever a construção de perfis estatísticos de classe robustos a valores
        extremos;
  \item reconhecer a estrutura do \ing{prompt} de um agente de diagnóstico e os
        critérios de validação de sua saída;
  \item explicar por que a tradução para linguagem intermediária é o que torna a camada
        de linguagem viável.
\end{itemize}
\end{objetivos}

\section{O problema de tradução}

Um modelo de linguagem não recebe séries temporais. Ele recebe texto.

A questão central deste capítulo é, portanto, uma questão de tradução: como converter um
segmento anômalo --- centenas de amostras, quatro sensores --- em uma descrição textual
que preserve a informação diagnóstica e seja compreensível pelo modelo.

Duas abordagens ingênuas falham de imediato. Despejar os números brutos no prompt gasta
milhares de \ing{tokens} em informação de baixíssima densidade, e o modelo não é bom em
extrair padrões de listas longas de floats. Descrever qualitativamente --- ``a pressão
subiu'' --- perde a magnitude, a velocidade e a relação entre sensores.

\begin{destaque}[A solução: uma linguagem intermediária]
Calcular deterministicamente um conjunto pequeno de métricas interpretáveis, e apresentar
ao modelo essas métricas em vez dos dados.

A escolha tem três consequências que valem ser explicitadas.

\textbf{Auditabilidade.} As métricas são calculadas por código verificável. Qualquer
afirmação do modelo pode ser confrontada com elas.

\textbf{Economia.} Treze números por sensor, contra centenas de amostras.

\textbf{Transferência de responsabilidade.} O trabalho de percepção fica com os métodos
numéricos --- que são bons nisso --- e o de interpretação com o modelo de linguagem.

O preço é que \emph{o que as métricas não capturam está perdido}. A escolha das métricas
é, portanto, a decisão de projeto mais consequente desta parte do livro.
\end{destaque}

\section{Pré-processamento}

A cadeia tem quatro etapas, nesta ordem.

\begin{enumerate}[leftmargin=*,itemsep=0.4em]

\item \textbf{Padronização por escore-$z$.} Cada sensor é centrado e escalado pela média
e pelo desvio padrão de sua referência. O conjunto usado para ajustar esses parâmetros
deve ser registrado para impedir vazamento e definir o que significa ``regime normal''.

\item \textbf{Normalização Min-Max para $[0,1]$.} Acomoda a entrada ao intervalo usado
pelo restante da cadeia. Se seus limites forem ajustados sobre as mesmas amostras da
etapa anterior, a composição é equivalente à normalização Min-Max direta; para ter
efeito adicional, as duas etapas precisam usar referências de ajuste distintas e
exclusivas do treinamento.

\item \textbf{Remoção de ruído por \ing{wavelet}.} Transformada de Daubechies db4, nível
3, com limiarização suave.

\item \textbf{Tratamento de valores ausentes.} Conforme o \cref{cap:dados}.

\end{enumerate}

\begin{destaque}[Por que db4, nível 3, limiar suave]
A escolha não é arbitrária, e cada componente responde a uma exigência do problema.

\textbf{Daubechies 4} tem suporte compacto e quatro momentos nulos --- suficiente para
representar bem transições abruptas, como o fechamento de uma válvula, sem introduzir os
artefatos oscilatórios que uma base mais suave produziria.

\textbf{Nível 3} remove três oitavas de alta frequência. Em sinais amostrados a cada
segundo, isso elimina ruído de instrumentação preservando dinâmicas de dezenas de
segundos.

\textbf{Limiarização suave} encolhe os coeficientes em vez de zerá-los, evitando as
descontinuidades que a limiarização dura introduz --- descontinuidades que seriam
interpretadas pelas métricas de \texttt{slope} e \texttt{rate} como eventos reais.
\end{destaque}

\subsection{Seleção de instâncias}

Uma decisão metodológica com consequência direta sobre a interpretação dos resultados:

\begin{itemize}[leftmargin=*,itemsep=0.3em]
  \item usam-se \textbf{apenas instâncias de poços reais} (\texttt{WELL});
  \item instâncias simuladas são acrescentadas somente quando uma classe tem menos de
        trinta instâncias reais;
  \item instâncias desenhadas à mão (\texttt{DRAWN}) são \textbf{sempre excluídas}.
\end{itemize}

A justificativa: o objetivo é avaliar a capacidade de interpretar fenômenos físicos
reais. Uma instância desenhada à mão contém a assinatura que um especialista
\emph{acredita} que o evento tem --- e um acerto sobre ela mediria concordância com a
intuição do anotador, não com a física.

\section{As treze métricas}

Para cada sensor, comparam-se dois trechos: o \termo{regime de base}{baseline}, anterior
ao evento, e o \termo{regime do evento}{event window}, delimitado pela segmentação do
\cref{cap:segmentacao}. Sejam $\bar{b}$ e $\sigma_b$ a média e o desvio padrão do
primeiro, e $\bar{e}$ e $\sigma_e$ os do segundo.

\subsection{Métricas de deslocamento}

\begin{align}
  \delta &= \bar{e} - \bar{b}, \label{eq:ag-delta} \\
  \text{variação percentual} &= 100 \, \frac{\delta}{|\bar{b}|}, \label{eq:ag-pct} \\
  z &= \frac{|\delta|}{\max(\sigma_b, 10^{-6})}. \label{eq:ag-z}
\end{align}

\begin{destaque}[As três métricas dizem coisas diferentes]
Uma pressão que sobe \SI{5}{\bar} sobre uma base de \SI{200}{\bar} com desvio padrão de
\SI{0,5}{\bar} tem $\delta = 5$, variação percentual de $2{,}5\,\%$ e $z = 10$.

O deslocamento absoluto é pequeno. A variação percentual é desprezível. E o escore-$z$
de $10$ indica um evento estatisticamente extraordinário.

É por isso que as três são fornecidas ao modelo. Cada uma responde a uma pergunta:
``quanto mudou?'', ``mudou muito em relação ao nível?'' e ``mudou muito em relação ao
que costuma variar?''. A última é quase sempre a diagnóstica, e é a que um humano
examinando o gráfico tende a subestimar.

O termo $\max(\sigma_b, 10^{-6})$ protege contra divisão por zero quando o sensor está
congelado --- situação frequente, como o \cref{cap:dados} discutiu.
\end{destaque}

\subsection{Métricas categóricas}

Três métricas convertem grandezas contínuas em categorias verbais, que o modelo de
linguagem manipula melhor que números.

\begin{center}
\small
\begin{tabular}{lll}
\toprule
\textbf{Métrica} & \textbf{Categoria} & \textbf{Critério} \\
\midrule
\multirow{3}{*}{\texttt{direction}}
& \ing{increase} & $\delta > 0{,}5\,\sigma_b$ \\
& \ing{decrease} & $\delta < -0{,}5\,\sigma_b$ \\
& \ing{stable}   & caso contrário \\
\midrule
\multirow{3}{*}{\texttt{magnitude}}
& \ing{small}    & $z < 1$ \\
& \ing{moderate} & $1 \leq z < 3$ \\
& \ing{large}    & $z \geq 3$ \\
\midrule
\multirow{3}{*}{\texttt{rate}}
& \ing{sudden}   & $|\text{inclinação}| > 0{,}01$ \\
& \ing{fast}     & $|\text{inclinação}| > 0{,}002$ \\
& \ing{gradual}  & caso contrário \\
\bottomrule
\end{tabular}
\end{center}

\begin{destaque}[Categorizar é uma decisão, não uma conveniência]
Fornecer ``a pressão aumentou de forma súbita e de grande magnitude'' em vez de
``$\delta = 4{,}7$, $z = 8{,}2$, inclinação $= 0{,}031$'' facilita o trabalho do modelo,
que foi treinado sobre texto e não sobre tabelas numéricas.

O custo é a perda de informação nas fronteiras: $z = 2{,}99$ e $z = 3{,}01$ recebem
categorias diferentes apesar de serem praticamente idênticos. Por isso as métricas
contínuas são fornecidas \emph{junto} com as categóricas, e não em vez delas. O modelo
recebe ambas e pode reconhecer casos de fronteira.
\end{destaque}

\subsection{Métricas de forma e variabilidade}

\begin{itemize}[leftmargin=*,itemsep=0.3em]
  \item \texttt{slope} --- inclinação da regressão linear sobre o segmento.
  \item \texttt{behavior} --- descrição qualitativa da forma da curva.
  \item \texttt{oscillation\_ratio} --- razão entre a energia oscilatória do evento e a
        da base. Discrimina golfadas de deslocamentos monotônicos.
  \item \texttt{noise\_ratio} $= \sigma_e / \sigma_b$ --- razão entre desvios padrão.
  \item \texttt{event\_range} --- amplitude do segmento.
  \item \texttt{baseline\_mean} e \texttt{event\_mean} --- os níveis absolutos.
\end{itemize}

\begin{destaque}[Por que fornecer os níveis absolutos]
As demais doze métricas são relativas, e há um bom motivo: elas precisam ser comparáveis
entre poços. Mas há informação diagnóstica no valor absoluto que o relativo apaga.

Uma temperatura de fundo de \SI{4}{\celsius} está próxima da condição de formação de
hidrato; uma de \SI{60}{\celsius} não está, e nenhuma variação relativa comunica isso.
Fornecer \texttt{baseline\_mean} e \texttt{event\_mean} permite ao modelo aplicar
conhecimento físico sobre faixas absolutas --- exatamente o tipo de raciocínio que os
métodos numéricos dos capítulos anteriores não fazem.
\end{destaque}

\subsection{Métricas entre sensores}

As treze métricas descrevem cada sensor isoladamente. Mas o \cref{cap:dominio} insistiu
que a informação diagnóstica está frequentemente na \emph{relação} entre sensores. Duas
métricas adicionais capturam isso:

\begin{itemize}[leftmargin=*,itemsep=0.3em]
  \item \textbf{relação P-PDG $\times$ P-TPT}, com três valores possíveis:
        \ing{same\_direction}, \ing{diverging}, ou
        \ing{one\_stable\_or\_both\_stable};
  \item \textbf{sensor dominante}, aquele com maior $|z|$.
\end{itemize}

\begin{exemplo}{Por que a divergência é diagnóstica}
Considere duas situações.

\emph{Primeira}: P-PDG e P-TPT caem juntas. O poço inteiro está despressurizando --- é
consistente com uma redução de produção, uma restrição a jusante, ou uma parada.

\emph{Segunda}: P-PDG sobe enquanto P-TPT cai. As duas pressões estão divergindo, o que
só é fisicamente possível se algo entre elas mudou --- tipicamente uma restrição ou o
fechamento de uma válvula na coluna.

O valor absoluto de cada pressão pode ser idêntico nos dois casos. É a relação que
distingue, e nenhuma métrica por sensor a captura.
\end{exemplo}

\section{Perfis de classe}

O modelo precisa saber contra o que comparar. Para cada classe conhecida constrói-se um
perfil estatístico a partir das instâncias de treinamento.

\subsection{Remoção de valores extremos em duas passagens}

Um perfil contaminado por uma instância atípica descreve mal a classe. A depuração usa o
critério do intervalo interquartil em duas passagens, com rigor crescente:

\begin{enumerate}[leftmargin=*,itemsep=0.2em]
  \item primeira passagem com $k = 3{,}0$ --- remove apenas extremos grosseiros;
  \item segunda passagem com $k = 1{,}5$ --- refina sobre a distribuição já limpa.
\end{enumerate}

\begin{destaque}[Por que duas passagens e não uma com $k = 1{,}5$]
Porque quartis são estimados a partir dos próprios dados. Um único valor muito extremo
desloca o terceiro quartil e, com ele, o limite superior --- fenômeno conhecido como
mascaramento. Aplicar $k = 1{,}5$ diretamente pode não remover o extremo que corrompeu a
estimativa.

Removendo primeiro os extremos grosseiros com $k = 3{,}0$, os quartis da segunda
passagem são estimados sobre dados mais limpos, e o critério mais rigoroso opera sobre
uma referência confiável. É um procedimento simples e frequentemente esquecido.
\end{destaque}

\subsection{Composição do perfil}

Cada perfil contém quatro elementos:

\begin{description}[leftmargin=0pt,style=nextline,itemsep=0.4em]
\item[Comportamento macro.] Padrão agregado ao longo de todo o evento.
\item[Comportamento local.] Padrão nos instantes iniciais --- útil para detecção precoce.
\item[Conhecimento de domínio.] Descrição física do mecanismo, redigida a partir da
literatura e da experiência operacional. É a única parte não derivada dos dados, e é
deliberadamente explícita: não se confia no que o modelo possa ter absorvido durante o
pré-treinamento.
\item[Características distintivas.] O que separa esta classe das mais parecidas com ela.
\end{description}

\begin{destaque}[O último elemento é o mais importante]
As três primeiras seções descrevem a classe. A quarta descreve suas \emph{fronteiras}.

Uma classe descrita isoladamente permite ao modelo reconhecer compatibilidade; descrita
em contraste com suas vizinhas, permite discriminar. Como o \cref{cap:resultados-agente}
mostrará, a maior parte dos erros do agente é entre classes fisicamente vizinhas --- e é
exatamente essa seção do perfil que ataca esse modo de falha.
\end{destaque}

\section{O prompt}

A estrutura final entregue ao modelo:

\begin{center}
\small
\begin{tabular}{ll}
\toprule
\textbf{Seção} & \textbf{Conteúdo} \\
\midrule
Papel        & especialista em análise de dados de poços submarinos \\
Escopo       & julgar apenas com base nas estatísticas fornecidas \\
Perfis       & as nove classes, com os quatro elementos acima \\
Evidência    & as treze métricas por sensor, mais as métricas entre sensores \\
Tarefa       & ordenar três candidatas; justificar; avaliar confiança \\
Novidade     & se nenhuma classe servir, nomear e descrever o fenômeno \\
Formato      & JSON com campos especificados \\
\bottomrule
\end{tabular}
\end{center}

A saída é validada automaticamente: JSON bem formado, exatamente três candidatas, classes
pertencentes ao conjunto declarado, confiança em $\{$\ing{High}, \ing{Medium},
\ing{Low}$\}$.

\begin{destaque}[O que este projeto \emph{não} inclui --- e por quê]
Não há ajuste fino do modelo, não há aprendizado com exemplos no prompt, não há múltiplas
passagens com agregação de votos, não há uso de ferramentas externas pelo agente.

Cada uma dessas técnicas provavelmente melhoraria os resultados. A omissão é deliberada:
o objetivo é estabelecer uma \textbf{linha de base honesta} para o que um modelo de
propósito geral, sem qualquer adaptação, consegue fazer nesta tarefa. Os números do
\cref{cap:resultados-agente} devem ser lidos como piso, não como teto.
\end{destaque}

\section{A arquitetura completa}

Vale, ao final desta descrição, situar a camada de agente no sistema inteiro.

\begin{enumerate}[leftmargin=*,itemsep=0.3em]
  \item A U-Net segmenta o sinal e delimita a região anômala (\cref{cap:segmentacao}).
  \item O classificador multiclasse produz a representação latente
        (\cref{cap:mahalanobis}).
  \item A distância de Mahalanobis decide se o segmento é conhecido ou novidade.
  \item As métricas deste capítulo são calculadas sobre o segmento.
  \item O agente recebe métricas e perfis, e produz hipóteses justificadas.
  \item O operador decide.
\end{enumerate}

\begin{destaque}[Cada camada depende da anterior]
A cadeia é sequencial, e isso significa que erros se propagam. Se a segmentação erra as
fronteiras --- como no poço B1, com IoU de $56\,\%$ ---, as métricas dos passos 4 são
calculadas sobre um segmento contaminado, e a interpretação do agente parte de premissas
erradas.

É a razão pela qual o \cref{cap:segmentacao} foi apresentado como pré-requisito da
explicabilidade, e não como um refinamento independente.
\end{destaque}

\resumocapitulo{%
A camada de agente resolve um problema de tradução: converter um segmento anômalo em
descrição textual que preserve informação diagnóstica. A solução é uma linguagem
intermediária de treze métricas por sensor --- deslocamento absoluto, variação
percentual, escore-$z$, inclinação, direção, magnitude, velocidade, comportamento, razão
de oscilação, razão de ruído, amplitude e os dois níveis médios --- acrescidas de duas
métricas entre sensores, entre as quais a relação entre P-PDG e P-TPT. Métricas
contínuas e categóricas são fornecidas simultaneamente, para que o modelo reconheça casos
de fronteira. O pré-processamento aplica padronização e normalização Min-Max com
referências de ajuste que precisam ser registradas, além de remoção de
ruído por \ing{wavelet} db4 nível 3 com limiarização suave, e usa apenas instâncias de
poços reais. Os perfis de classe, depurados por remoção de extremos em duas passagens com
$k = 3{,}0$ e depois $k = 1{,}5$, contêm comportamento macro, comportamento local,
conhecimento de domínio explícito e --- o elemento decisivo --- as características que
distinguem a classe de suas vizinhas. O agente recebe perfis e evidência, e produz três
candidatas ordenadas com justificativa e confiança, em JSON validado. Não há ajuste fino
nem agregação de múltiplas passagens: os resultados do próximo capítulo são um piso.}

\begin{exercicios}
\item \textbf{Projeto de extensão.} Implemente as treze métricas e aplique-as a
      instâncias de cada classe do 3W.
      Construa uma matriz de correlação entre elas: quais são redundantes? Proponha um
      subconjunto mínimo que preserve a capacidade discriminativa.

\item O escore-$z$ usa $\max(\sigma_b, 10^{-6})$. Analise o comportamento da métrica
      quando o sensor de base está congelado ($\sigma_b = 0$) mas o evento apresenta
      variação. O valor resultante é informativo ou é um artefato? Proponha um
      tratamento alternativo.

\item As categorias de \texttt{magnitude} usam fronteiras em $z = 1$ e $z = 3$.
      Justifique essas escolhas sob a hipótese de normalidade e discuta o que acontece
      quando ela é violada. Proponha fronteiras baseadas em quantis empíricos e compare.

\item Implemente a remoção de extremos em duas passagens e compare com uma única passagem
      com $k = 1{,}5$. Construa um conjunto sintético em que o mascaramento ocorra e
      quantifique a diferença nos perfis resultantes.

\item Projete uma métrica adicional entre sensores que capture a defasagem temporal
      entre a resposta de P-PDG e a de P-TPT. Que classes ela ajudaria a discriminar?
      Implemente-a e verifique empiricamente.

\item \textbf{Atividade com serviço externo.} Escreva o prompt completo descrito neste
      capítulo, registre modelo, revisão, data e parâmetros, e submeta o mesmo segmento
      vinte vezes com temperatura $0{,}2$. Meça a variabilidade da saída. Repita com
      temperatura $0{,}0$ e $0{,}7$ e discuta o compromisso entre consistência e
      diversidade de hipóteses.

\item Argumenta-se que instâncias desenhadas à mão devem ser excluídas porque medem
      concordância com a intuição do anotador. Construa o argumento contrário: em que
      circunstâncias avaliar sobre instâncias sintéticas seria legítimo e informativo?
\end{exercicios}

\begin{leituras}
\item \citet{Lopes2026LLM} --- o trabalho que originou este capítulo, com o prompt
      completo.
\item \citet{VARGAS2019} e \citet{VazVargas2026_3w2} --- a descrição do 3W e das
      proveniências das instâncias.
\item \citet{Ma2024wellrecords} e \citet{Wei2024leakage} --- aplicações de modelos de
      linguagem a registros e a diagnóstico no setor.
\item \citet{Belay2026agenticAD} --- detecção de anomalias com arquiteturas de agente.
\item \citet{Zhang2026llmrules} --- extração de regras de domínio com modelos de
      linguagem.
\end{leituras}
